% N1 supplementary material (Journal of Process Control). Build from this folder like main.tex:
% pdflatex supplement (twice). Tables are read from the sibling IPIS checkout, written by
% IPIS scripts/paper_n1/supplement.py from the Phase 1 evidence.
\documentclass[10pt,a4paper]{article}
\usepackage[margin=2cm]{geometry}
\usepackage{mathptmx}
\usepackage{amsmath}
\usepackage{booktabs}
\usepackage{array}
\usepackage{url}
\usepackage{natbib}
\usepackage[labelfont=bf]{caption}
\newcommand{\figdir}{../../IPIS/docs/paper/figures/n1}
\renewcommand{\thetable}{S\arabic{table}}
\renewcommand{\thesection}{S\arabic{section}}
\setlength{\parskip}{3pt}
\setlength{\parindent}{0pt}

\title{Supplementary material for ``When the laboratory is late: dead-time analysis and
compensation of conformal soft-sensor intervals''}
\author{Bien Don Busico\\
\small University of the Philippines Diliman, Quezon City, Philippines\\
\small \texttt{bienbusico@gmail.com}}
\date{}

\begin{document}
\maketitle

This supplement holds what the article refers to: every cell and contrast of the pre-registered
study (Section~\ref{sec:files} and Tables~\ref{tab:T1} to~\ref{tab:baselines}), the run and its
checks (Section~\ref{sec:run}), the Tennessee Eastman simulator's random numbers
(Section~\ref{sec:rng}), the baseline configurations (Section~\ref{sec:baselines}), and the design
checks behind the rules (Section~\ref{sec:design}). Notation and abbreviations follow the article:
$\theta$ is the dead time in calibrator updates, $U$ the share of issued intervals that are
infinite or wider than the range of $y$ over the calibration data, $\varepsilon = 1\,\%$, and the
replay targets $\varepsilon/2$ for a budget of $2\varepsilon$. Every number is read from the records
of the run (stages A and B; IPIS repository, commit \texttt{bd6454d},
\texttt{docs/paper/evidence/}) or from the pre-registration and its pilot evidence, as stated. The
three Tennessee Eastman operating points are labelled mode 1 to 3 in the records and tables.

\section{Data files}\label{sec:files}

Three files accompany this document, written from the stage B records by
\texttt{scripts/paper\_n1/supplement.py} in the IPIS repository and archived with the evidence.
\begin{itemize}\setlength{\itemsep}{0pt}
\item \texttt{n1\_supp\_cells.csv}: every loop cell of every unit (978 cells). Key fields: unit
(T1, T1\_gb, T2, T3, debutanizer), $\theta$, compensation (raw or comp), loop (none, sequential,
projected, phase, smith, decaying), step-size policy, pairing (stored or arrival), the number of
runs, and the step size and tuning feasibility per operating point. Metrics, each with its mean
and 95\,\% interval: $U$, the infinite share, coverage, the median, 90th and 99th percentiles of
the finite widths, F7's width (the median over all issued intervals), the worst rolling coverage
over 100 intervals, sd($\alpha$), and the IACT, lag-1 autocorrelation and mean run length of the
misses. Tennessee Eastman units: means over the 15 runs with $t$ intervals where the stage B
summary holds them (run means without an interval otherwise). Debutanizer: the single stream,
with circular block bootstrap intervals for $U$, the infinite share and coverage.
\item \texttt{n1\_supp\_contrasts.csv}: every paired contrast (18\,466 rows): compensated minus
uncompensated, each $\theta$ minus $\theta = 0$, arrival minus stored pairing, loop on minus loop
off, each loop minus the sequential loop, and each policy minus replay to $\varepsilon/2$; mean and
95\,\% interval (Tennessee Eastman: $t$ over the runs; debutanizer: block bootstrap over all
evaluation rows).
\item \texttt{n1\_supp\_baselines.csv}: the four baseline configurations at every $\theta$ of the
units that ran them (T1 and the debutanizer).
\end{itemize}

\section{The run and its checks}\label{sec:run}

The protocol was frozen at tag \texttt{n1-phase1-freeze} (29 October 2026). The run, on
2 November 2026 (UTC), had three commits.
\begin{itemize}\setlength{\itemsep}{0pt}
\item Runner, \texttt{5ed6f0c}: the files of the checked runner, byte for byte.
\item Stage A, \texttt{133b0f3}: training and calibration rows only (613~s on six workers). It
fitted the point models and compensators, tuned every step size, configured the baselines and
recorded the predicted width ratios. Its records were committed and pushed before stage B read an evaluation
row.
\item Stage B and the verdicts, \texttt{bd6454d} (1822~s). Stage B recomputed stage A's settings,
streams, file digests, point models, arrays (by SHA-256), compensator orders and coefficients,
predicted ratios, static quantiles, ranges of $y$ and conformal PID replays, and passed all 1359
checks exactly before it evaluated.
\end{itemize}
A recomputation of the verdicts from the committed stage A and stage B files on a second platform
(Linux) reproduced them: of 4772 values, 236 differ in the last digits, with a largest relative
difference of $1.3 \times 10^{-9}$. In T1's stage B, statsmodels warned twice of an overflow inside its bounded transform of the Theta
model's smoothing weight (conformal PID's Theta scorecaster); the transform then returns its
bound, every check was exact and the cell's results are finite.

One deviation from the protocol was logged before any result: arm T3, which updates once per
analyzer value, seeds its window with 40 labels instead of 200, because its 40 calibration hours
hold only 160 labels. No rule is decided on T3.

Stage A flagged nine tunings as infeasible, all at T2's first operating point without
compensation: replay to $\varepsilon/2$ at $\theta = 4$, and to $\varepsilon$ and $\varepsilon/2$
at $\theta = 14$, for the sequential, projected and Smith loops. There $U$ exceeded its target
already at the grid's smallest step (0.75 to 0.76\,\% at $\theta = 4$ against 0.5\,\%; 1.72 to
1.73\,\% at $\theta = 14$ against targets of 1 and 0.5\,\%), so those cells ran at 0.001, as the
protocol fixes. No tuning was infeasible at T1, T1\_gb, T3 or the debutanizer, and every conformal PID tuning met its
selection criterion. The ridge penalty chosen by the one-standard-error rule was the top of its
grid, 1000, in 10 of the 15 runs.

\section{The Tennessee Eastman simulator's random numbers}\label{sec:rng}

The simulator's generator, TESUB7 in \texttt{teprob.f}, computes
$G \leftarrow 9228907\,G \bmod 2^{32}$ in double precision. The product exceeds $2^{53}$, so IEEE
double arithmetic rounds it; that rounding is part of the algorithm was recorded by
\citet{kitchin2026s}. We add what follows from it. The state gains trailing zero bits until
$G = 2^k m$ with $m$ odd, after which the recursion is exact with period $2^{30-k}$ draws. A
counting build gives 952\,326 draws per simulated hour. Over the source's own seed and the 57
seeds listed in \texttt{teprob.f}, the settled $k$ and the periods are:

\begin{center}\footnotesize
\begin{tabular}{@{}lrrrrrrrr@{}}
\toprule
Settled $k$ & 3 & 4 & 5 & 6 & 7 & 8 & 9 & 11 \\
Seeds & 17 & 24 & 8 & 2 & 4 & 1 & 1 & 1 \\
Period of the random stream (h) & 141 & 70.5 & 35.2 & 17.6 & 8.8 & 4.4 & 2.2 & 0.55 \\
\bottomrule
\end{tabular}
\end{center}

So 41 of the 58 seeds repeat their random stream within 100~h; with $k = 4$, hours 70 to 100 of a
run would reuse the random numbers of hours 0 to 30. For each $k$ the odd $m$ form two cycles (the
multiplier is 3 mod 8), so seeds with the same $k$ and cycle draw the same numbers at different
phases and do not give independent runs. These periods hold for IEEE double arithmetic, as with
gfortran~13 on x86-64.

For the study, each run drew instead from its own stream of MRG32k3a \citep{lecuyer1999s}, exact
in double precision: run $j$ starts at stream $j$, $j \times 2^{127}$ draws from the RngStreams
seed, and a 100-h run uses about $10^{8}$ draws, so no two runs share a random number. The Fortran
code reproduces a Python reference draw for draw (200\,000 draws), and its first output and stream
jumps agree with an independent implementation. The plant model, the disturbances and the
controllers are unchanged. Streams were dealt to the three operating points in rotation, and each
took the first five whose 100-h run did not shut down: 5 of 26 candidates survived at the first
operating point, 14 of 25 at the second and 17 of 25 at the third.

\section{Baseline configurations}\label{sec:baselines}

The settings follow each paper and, where a paper is silent, the authors' code: conformal PID
control (\citealp{angelopoulos2023pids}; code \texttt{aangelopoulos/conformal-time-series}, commit
\texttt{b729c3f}) and SPCI (\citealp{xu2023spcis}; code \texttt{hamrel-cxu/SPCI-code}, commit
\texttt{d598dd3}). Where code contradicts its paper, or would use a label before it arrives, the
paper governs. Appendix A of the pre-registration gives every setting with its source, and every
departure with its reason. Among them: conformal PID starts from the loop's seed rows, its
learning rate uses the range of the recent scores (as the code does) rather than their maximum
(as the paper does), and its quantile-regression scorecaster has no penalty; SPCI's window, which
the paper leaves open, is scaled from the authors' setting on their electricity data (300 for
2755 training residuals) to the calibration length.

\emph{Conformal PID.} Signed residual scores, each side tracked at $\alpha/2 = 0.05$. The update
for sample $t+1$ uses the coverage error of the newest labelled sample, $\theta + 1$ steps back.
Learning rate $m$ times the range of the 200 latest labelled scores; integrator
$K_I \tan(x \log n /(C_{\mathrm{sat}} n))$. Scorecasters: (a) a linear quantile regression at level
0.95 on the three latest labelled scores at the horizon; (b) the Theta model, refitted at every
step; (c) the compensated loop's own residual forecast. Tuning: 114 configurations ($m$ on 19
values from 0.001 to 1; $C_{\mathrm{sat}}$ on three values; $K_I$ at $B'$ and $10B'$, with $B'$ the
largest absolute calibration residual), replayed on the pooled calibration segments from the same
seeded start as the loop. Among configurations with replayed $U \le \varepsilon/2$ and coverage
within 0.02 of 0.90, the one with the smallest replayed median width, scaled to 0.90 coverage, was
taken.

\emph{SPCI.} The point model's residuals in place of EnbPI's leave-one-out ensemble; one quantile
regression forest for horizon $\theta + 1$ (10 trees, maximum depth 2), trained on windows of
$w$ known residuals with the residual $\theta + 1$ steps on as the target and refitted at every
step on the latest $T_{\mathrm{cal}}$ known residuals; $w = \operatorname{round}(300\,T_{\mathrm{cal}}
/2755)$, 87 for T1 and 181 for the debutanizer; $\hat\beta$ from five values minimising the
predicted width.

F7 scales each baseline's width to the loop's coverage in each run by the ratio of Gaussian
quantiles. That scaling errs towards whichever method covers less: in a dry run before the study,
the closest competitor's coverage-matched difference was $-2.8$\,\% under the Gaussian scaling
and $-3.7$ to $-4.5$\,\% under $t_5$, $t_3$ and Laplace tails, against F7's 5\,\% margin. The
comparison scaled by the loop's own score quantiles is therefore reported beside it: +1.7,
+17.9, $-1.3$ and +2.1\,\% for the four baselines in the order of the article's Table~4, against
+1.6, +17.7, $-1.5$ and +1.6\,\% under the Gaussian scaling.

\section{Design checks}\label{sec:design}

\emph{Power.} On synthetic AR(1) residuals analysed with the study's code (200 studies at each
persistence), F2's test detected the delay effect at the delay-blind step with probability 0.91
($\phi = 0.5$) and 0.99 ($\phi = 0.9$) in the Tennessee Eastman design (15 runs of 800 calibration
and 600 evaluation rows, $\theta = 4$), but only 0.12 and 0.095 on a single stream of the
debutanizer's length ($\phi = 0.965$ and 0.970, $\theta = 3$). Tuned to $\varepsilon$, the
uncompensated loop's 15-run mean $U$ averaged 0.59 and 0.93\,\% but exceeded $\varepsilon$ in 21
and 38\,\% of studies; tuned to $\varepsilon/2$ it exceeded $2\varepsilon$ in 0 and 4.0\,\%
(compensated: 0 and 2.5\,\%), which set F5's tuning target and budget. The mean over 15 runs of
the width-ratio error lay within $\pm 10\,\%$ in every study at both persistences.

\emph{The white-noise model.} On synthetic residuals (absolute values of a Gaussian AR(1) series,
window 200, five seeds), the article's Eq.~(7) predicted the standard deviation of the level to
within $-6$ to $+5$\,\% for steps of at least 0.02 when the series had no serial correlation
($\phi = 0$), but 22 to 38\,\% too high at a step of 0.005, because the sliding window's quantile
adds feedback of its own. At $\phi = 0.9$ and steps of at least 0.02 it was 47 to 63\,\% too low.

\emph{Width ratio and residual shape.} Eq.~(9) of the article is the ratio of standard deviations;
it equals the width ratio when the compensated and uncompensated scores share one standardised
shape. For AR(1) residuals with $\phi = 0.9$ and uniform or $t_5$ innovations, simulated over 2
million steps, the width ratio at a one-step horizon was 4.2 to 5.5\,\% below the ratio of
standard deviations.

\emph{F8.} F8 is tested only where Eq.~(9) predicts a gain (mean predicted ratio at most 0.95).
On Gaussian AR(1) residuals without a label hold it fired in 1.4\,\% of tested studies at
$\phi = 0.8$ and in none at 0.85 to 0.95; with a fifth of the residual variance held for five rows,
as T1's analyzer holds its reading, it fired in at most 0.5\,\%.

\section{Tables}

Tables~\ref{tab:T1} to~\ref{tab:deb} give, for each unit, the sequential loop with stored pairing
under every step-size policy, and the loop off (split conformal), without and with compensation.
Tennessee Eastman: means over the 15 runs; the step is given per operating point (1/2/3). Widths
are full widths in the units of $y$ (mol\,\% for the Tennessee Eastman process), the run median of
the finite widths; ``worst rolling'' is the minimum coverage over 100 consecutive intervals; IACT
is that of the miss indicator. The intervals and every other cell are in the data files.

\begin{table}[p]
\centering\footnotesize
\caption{T1 (3-min updates; ridge regression): every policy of the sequential loop.}
\label{tab:T1}
\input{\figdir/n1_supp_primary_T1.tex}
\end{table}

\begin{table}[p]
\centering\footnotesize
\caption{T1\_gb (3-min updates; gradient boosting): every policy of the sequential loop.}
\label{tab:T1gb}
\input{\figdir/n1_supp_primary_T1_gb.tex}
\end{table}

\begin{table}[p]
\centering\footnotesize
\caption{T2 (1-min updates; ridge regression): every policy of the sequential loop.}
\label{tab:T2}
\input{\figdir/n1_supp_primary_T2.tex}
\end{table}

\begin{table}[p]
\centering\footnotesize
\caption{T3 (one update per analyzer value; window of 40 labels; descriptive): every policy of
the sequential loop.}
\label{tab:T3}
\input{\figdir/n1_supp_primary_T3.tex}
\end{table}

\begin{table}[p]
\centering\footnotesize
\caption{Debutanizer (one stream, 719 evaluation rows): every policy of the sequential loop. The
step is the single tuned value.}
\label{tab:deb}
\input{\figdir/n1_supp_primary_debutanizer.tex}
\end{table}

\begin{table}[p]
\centering\footnotesize
\caption{The comparison loops at T1, $\theta = 4$, under every step-size policy: means over the 15
runs, uncompensated, with the compensated loop's $U$ in the last column. Projected: the level
clipped to $[0.01, 0.5]$; interleaved: $\theta + 1$ delay-free chains (El Halabi and Brandt);
Smith form: the miss rate expected at the stored level replaced by that at the current level;
decaying: steps $c\,n^{-0.6}$.}
\label{tab:loops}
\input{\figdir/n1_supp_loops.tex}
\end{table}

\begin{table}[p]
\centering\footnotesize
\caption{The four baseline configurations beside the compensated loop tuned by replay to
$\varepsilon/2$, at every $\theta$ of T1 (means over the 15 runs) and of the debutanizer.
F7 width: the run median over all issued intervals, infinite ones included. F7 is decided at
T1, $\theta = 4$ only; the other rows are descriptive.}
\label{tab:baselines}
\input{\figdir/n1_supp_baselines.tex}
\end{table}

\clearpage
\begin{thebibliography}{9}
\setlength{\itemsep}{0pt}
\bibitem[Angelopoulos et al.(2023)]{angelopoulos2023pids} A. Angelopoulos, E. Cand\`es,
R. J. Tibshirani, Conformal PID control for time series prediction, in: Advances in Neural
Information Processing Systems (NeurIPS 2023), vol. 36, 2023, pp. 23047--23074.
\bibitem[Kitchin(2026)]{kitchin2026s} J. Kitchin, tep-rust: pure-Rust Tennessee Eastman process
simulator, GitHub repository, \url{https://github.com/jkitchin/tep-rust}, development log, entry
of 10 August 2026; commit 0735ad3.
\bibitem[L'Ecuyer(1999)]{lecuyer1999s} P. L'Ecuyer, Good parameters and implementations for
combined multiple recursive random number generators, Operations Research 47 (1) (1999)
159--164, \url{https://doi.org/10.1287/opre.47.1.159}.
\bibitem[Xu and Xie(2023)]{xu2023spcis} C. Xu, Y. Xie, Sequential predictive conformal inference
for time series, in: Proceedings of the 40th International Conference on Machine Learning, PMLR
202, 2023, pp. 38707--38727.
\end{thebibliography}

\end{document}
